\documentclass[12pt,a4paper]{article}
\usepackage{ugmath}
\usepackage{float}
\usepackage{placeins}
\usepackage [spanish]{babel}
\usepackage[labelfont=bf,labelsep=space]{caption}
\newcommand{\paren}[1]{\left( #1 \right)}
\newcommand{\alumno}{Ricardo León Martínez}
\newcommand{\materia}{Calculo Diferencia e Integral II}
\newcommand{\profesor}{Fernando Nuñez Medina}
\newcommand{\tarea}{Tarea 4}
\newcommand{\fecha}{20/2/2026}

\begin{document}
    \begin{center}
    {\large\textbf{UNIVERSIDAD DE GUANAJUATO}}\\[0.3cm]
    {\normalsize\textbf{DIVISIÓN DE CIENCIAS NATURALES Y EXACTAS}}\\
    {\normalsize\textbf{CAMPUS GUANAJUATO}}\\[1cm]

    {\Large\textbf{\tarea\ (\materia)}}\\[1cm]
    \end{center}

    \textbf{Nombre:} \alumno \hfill 
    \textbf{Fecha:} \fecha \hfill 
    \textbf{Calificación:} \rule{3cm}{0.4pt} \\[0.3cm]

    \begin{mdframed}[style=mdbluebox,frametitle={Ejercicio 1}]
        Prueba la proposición 12. \textbf{Sugerencia:} Toda partición $P=\{t_{0},\dots,t_{n}\}$ del intervalo $[a,b]$ genera una partición
        $Q=\{t_{0}-c,\dots,t_{n}-c\}$ del intervalo $[a-c,b-c]$ y viceversa.
    \end{mdframed}

    \begin{mdframed}[style=mdbluebox,frametitle={Ejercicio 2}]
        Calcula las siguientes integrales:
        \begin{enumerate}
            \item[(a)] $\displaystyle\int\frac{x}{x^{2}+1}dx$.
            \item[(b)] $\displaystyle\int\sin^{2}(x)\cos(x)dx$.
            \item[(c)] $\displaystyle\int\frac{x}{\sqrt{4-x^{2}}}dx$.
            \item[(d)] $\displaystyle\int\frac{x}{\sqrt{9+x^{2}}}dx$.
            \item[(e)] $\displaystyle\int\frac{x}{\sqrt{x^{2}-2}}dx$.    
        \end{enumerate}
    \end{mdframed}

    \begin{mdframed}[style=mdbluebox,frametitle={Ejercicio 3}]
        Calcula la integral
        \begin{equation*}
            \int\frac{x^{2}+5x-14}{(x+4)(x^{2}+2)}dx.
        \end{equation*}
    \end{mdframed}

    \begin{mdframed}[style=mdbluebox,frametitle={Ejercicio 4}]
        Considera las funciones $f:(0,1]\to\mathbb{R}$,$g:(0,1]\to\mathbb{R}$ y $h:(0,1]\to\mathbb{R}$ dadas por
        \begin{equation*}
            f(x)=\frac{1}{\sqrt{x}},\quad g(x)=\frac{1}{x}\text{ y }h(x)=\frac{1}{x^{2}}
        \end{equation*}
        \begin{enumerate}
            \item[(a)] Dibuja las graficas de $f,g$ y $h$.
            \item[(b)] Calcula el área bajó la gráfica de cada una de las funciones $f,g$ y $h$.  
        \end{enumerate}
    \end{mdframed}

    \begin{mdframed}[style=mdbluebox,frametitle={Ejercicio 5}]
        Calcula lo siguiente:
        \begin{enumerate}
            \item[(a)] El área entre las gráficas de las funciones $f(x)=x^{2}$ y $g(x)=1$ sobre el intevervalo
            $[-1,1]$.
            \item[(b)] El área entre las gráficas de las funciones $f(x)=\sin(x)$ y $g(x)=\cos(x)$ sobre el intervalo
            $[0,\pi]$.
        \end{enumerate}
    \end{mdframed}

    \begin{mdframed}[style=mdbluebox,frametitle={Ejercicio 6}]
        Muestra que la longitud de la circunferencia con centro en el origen  y radio $r$ es $2\pi r$.
    \end{mdframed}

    \begin{mdframed}[style=mdbluebox,frametitle={Ejercicio 7}]
        Utiliza la circunferencia trigonometrica para calcular, si es que existe, el valor de las funciones trigonometricas del angulo
        $\theta$ de cada uno de los incisos siguientes (haz un dibujo).
        \begin{multicols}{2}
            \begin{enumerate}
                \item[(a)]  $\theta=\pi.$
                \item[(b)]  $\theta=3\pi/2.$
                \item[(c)]  $\theta=2\pi.$
                \item[(d)]  $\theta=-\pi/2.$
            \end{enumerate}
        \end{multicols}
    \end{mdframed}

    \begin{mdframed}[style=mdbluebox,frametitle={Ejercicio 8}]
        Prueba los incisos (c) al (f) de la proposición 23.
    \end{mdframed}
\end{document}